\begin{table}[tp]
\centering
\caption{Finite-variance generalized-smooth convex optimization: \href{https://github.com/k3nfalt/longtails/blob/main/LongTails/Setups/FiniteVarianceGeneralizedSmoothAndConvex.lean}{FiniteVarianceGeneralizedSmoothAndConvex}.}\label{tab:FiniteVarianceGeneralizedSmoothAndConvex}
\parbox{\linewidth}{\footnotesize Convex $f$ on a separable real Hilbert space has differentiable gradient and $\|\nabla^2 f(x)\|\le L_0+L_1\|\nabla f(x)\|$. A minimizer $x^\star$ satisfies $\|x_0-x^\star\|\le R$ and $f(x_0)-f(x^\star)\le\Delta$. IID samples with marginal $p$ and joint law $P$ feed a jointly measurable oracle $g$, with integrable $g(x,\cdot)$ and squared error, $\int g(x,\omega)\,dp=\nabla f(x)$ and $\int\|g(x,\omega)-\nabla f(x)\|^2\,dp\le\sigma^2$. Known parameters satisfy $L_0,R,\Delta,\varepsilon>0$, $L_1,\sigma\ge0$, $0<\delta<1$. Target: measurable final output $\hat x$ with $P\{f(\hat x)-f(x^\star)>\varepsilon\}\le\delta$. Costs count all stochastic-gradient samples, including pilots, corrections, and padding.}\par\smallskip
\begin{threeparttable}
\footnotesize
\setlength{\tabcolsep}{4pt}\renewcommand{\arraystretch}{1.25}
\begin{tabularx}{\linewidth}{@{}>{\raggedright\arraybackslash}p{0.20\linewidth}>{\raggedright\arraybackslash}p{0.32\linewidth}>{\raggedright\arraybackslash}p{0.25\linewidth}>{\raggedright\arraybackslash}X@{}}
\toprule
\textbf{Result} & \textbf{Oracle complexity} & \textbf{Novelty} & \textbf{Limitations} \\
\midrule
\textbf{(1)} \href{https://github.com/k3nfalt/longtails/blob/main/LongTails/Theorems/AdaptiveRobustGeneralizedAcceleratedConvex.lean}{Adaptive\allowbreak Robust\allowbreak Generalized\allowbreak Accelerated\allowbreak Convex} \textbf{(Ours)} $(\star\star\star)$ \href{https://github.com/k3nfalt/longtails/blob/main/scripts/reviews/AdaptiveRobustGeneralizedAcceleratedConvex.md}{(review)}
& $\begin{aligned}[t]
\mathcal{O}\Bigl(& (W+v\ell)\Bigl[1+\frac{\log(eM/\delta)}{\log(e+W\sqrt M/(v^2\ell^2))}\Bigr]\\
&+v^2\ell+v\ell\log(eM/\delta)\Bigr).
\end{aligned}$
\par For fixed $c>0$, $\delta\ge(e+W)^{-c}$ implies $\mathcal{O}_c(W+v^2\ell)$.
& & \\
\midrule
\textbf{(2)} \href{https://github.com/k3nfalt/longtails/blob/main/LongTails/Theorems/FiniteVarianceGeneralizedAcceleratedConvex.lean}{Finite\allowbreak Variance\allowbreak Generalized\allowbreak Accelerated\allowbreak Convex} \textbf{(Ours)} $(\star\star)$ \href{https://github.com/k3nfalt/longtails/blob/main/scripts/reviews/FiniteVarianceGeneralizedAcceleratedConvex.md}{(review)}
& $\mathcal{O}\!\left((W+v^2\ell^2)\log\frac{e(1+T+q+\lambda)}{\delta}\right)$ & & \\
\midrule[1.2pt]
\textbf{(3)} RCSAG \citep{fvLobanov2026v1}, Thm.~9, App.~C; derived bound on $\mathbb R^d$, $p=2$, constant variance, $\Delta>\varepsilon$.
& $\mathcal{O}\!\left((\lambda+\sqrt T+qJ+v^2)\log\frac{e(1+\lambda)}{\delta}\right)$
& \textbf{(1)}: removes the $\ell$ factor on $\sqrt T$ when $\delta\ge(e+W)^{-c}$ and optimization dominates.
\par \textbf{(1)}, \textbf{(2)}: for $T=1$, $v\le1$, $\Delta=4\varepsilon$, $\delta=1/2$, $q\to\infty$, cost $\mathcal{O}(q)$ becomes $\mathcal{O}(q^{2/3})$, $\mathcal{O}(q^{2/3}\log q)$, respectively. Extend to separable Hilbert spaces.
& \textbf{(1)}, \textbf{(2)}: require a Hessian bound and finite variance; \textbf{(3)} permits first-order local smoothness, $p<2$, gradient/gap-dependent noise.
\par \textbf{(1)}: worse for $T=v=1$, $q=0$, $\Delta=2\varepsilon$, $\delta\to0$: $\mathcal{O}(\ell^2)$ versus $\mathcal{O}(\ell)$.
\par \textbf{(2)}: extra noise and optimization logarithms. \\
\midrule
\textbf{(4)} Clipped accelerated mirror descent \citep{fvNguyen2023v2}, Thm.~4.1; special case $L_1=0$, $p=2$, Euclidean geometry.
& $\mathcal{O}\!\left(1+(\sqrt T+v^2)\ell\right)$
& \textbf{(1)}: $\mathcal{O}(\sqrt T)$ versus $\mathcal{O}(\sqrt T\log T)$ for $T\to\infty$, $\delta=T^{-1}$, $v^2\ell\le\sqrt T$, using $\Delta\le L_0R^2/2$. No gain at constant $\delta$. \textbf{(1)}, \textbf{(2)} extend to $L_1>0$.
& \textbf{(1)}, \textbf{(2)} need $\Delta$ and differentiable gradients; \textbf{(4)} allows $p<2$ and mirror geometries. \textbf{(1)} has the tiny-$\delta$ loss shown against \textbf{(3)}. \\
\midrule
\textbf{(5)} \citep{fvTyurin2026v2}, Thm.~4.3; special case $\sigma=0$, exact gradients, $\Delta>\varepsilon$.
& $\begin{aligned}[t]\mathcal{O}\Bigl(&1+\sqrt T+q\bigl[1+\log_+\min\{\tfrac{L_1^2\Delta}{L_0},\tfrac{\Delta}{\varepsilon}\}\bigr]\\
&+\log_+\tfrac{\Delta}{L_0R^2}\Bigr)\end{aligned}$
& \textbf{(2)}: at $T=1$, $\Delta=4\varepsilon$, $\delta=1/2$, $q\to\infty$, $\mathcal{O}(q^{2/3}\log q)$ versus $\mathcal{O}(q)$. \textbf{(1)}, \textbf{(2)} extend to finite-variance stochastic gradients.
& \textbf{(2)} pays $\log(e(1+T+q+\lambda)/\delta)$ on $\sqrt T$ when it dominates, even with exact gradients. \\
\midrule
\textbf{(6)} \href{https://github.com/k3nfalt/longtails/blob/main/LongTails/Theorems/FiniteVarianceGeneralizedRadiusAcceleratedConvex.lean}{Finite\allowbreak Variance\allowbreak Generalized\allowbreak Radius\allowbreak Accelerated\allowbreak Convex} \textbf{(Ours)} $(\star\star\star)$ \href{https://github.com/k3nfalt/longtails/blob/main/scripts/reviews/FiniteVarianceGeneralizedRadiusAcceleratedConvex.md}{(review)}
& $\mathcal{O}\!\left((1+q+\sqrt T+v^2)H^3\right)$,
\par $H=\log\!\dfrac{e(2+T+q+v)}{\delta}$.
& & \\
\midrule[1.2pt]
\textbf{(7)} Clip-SSTM \citep{fvRadiusElKadi2026v2}, Thm.~2, Remark~1; derived bound on $\mathbb R^d$, $L_1>0$, $\alpha=2$, supplied gap.
& $\begin{aligned}[t]\mathcal{O}\Bigl(\bigl(& (1+q)[D+\log((1+q)H_D)]\\
&+\sqrt T+v^2\bigr)H_D\Bigr),\end{aligned}$
\par $D=\log(2+\Delta/(L_0R^2))$,
$H_D=\log\!\dfrac{e(2+T+q+v+D)}{\delta}$.
& \textbf{(6)}: curvature $qD\to q$ up to logs; gain when large $D$ dominates logarithmic overhead and unchanged terms. Updates and budget need no $\Delta$.
& \textbf{(6)}: extra logarithmic overhead and larger leading constant; covers only $\alpha=2$, whereas \textbf{(7)} allows $1<\alpha<2$. \\
\midrule
\textbf{(8)} Clip-SSTM \citep{fvRadiusElKadi2026v2}, Thm.~2, Remark~2; derived radius-only bound on $\mathbb R^d$, $L_1>0$, $\alpha=2$.
& $\begin{aligned}[t]\mathcal{O}\Bigl(\bigl(& (1+q)[q+\log((1+q)H)]\\
&+\sqrt T+v^2\bigr)H\Bigr).
\end{aligned}$
& \textbf{(6)}: curvature $q^2\to q$ up to logs. For $T=v=1$, $\delta=1/2$, $q\to\infty$: $\mathcal{O}(q\log^3q)$ versus $\mathcal{O}(q^2\log q)$. No total gain implied when unchanged terms dominate.
& \textbf{(6)}: extra logarithmic overhead and larger leading constant; excludes the $\alpha<2$ noise allowed by \textbf{(8)}. \\
\midrule
\textbf{(9)} Earlier repository bound \citep{fvRadiusEarlierRepository}; comparison with \textbf{(2)} on the identical setup.
& $\mathcal{O}\!\left((W+v^2\ell^2)\log\frac{e(1+T+q+\lambda)}{\delta}\right)$
& \textbf{(6)}: no $\Delta$ needed. For $\lambda\asymp q\to\infty$, $T=v=1$, $\delta=1/2$, admissible cosh instances give $\mathcal{O}(q\log^3q)$ versus $\mathcal{O}(q^{4/3}\log q)$.
& \textbf{(6)}: worse for small known gaps; e.g. fixed $\lambda,T,v,\delta$ gives $\mathcal{O}(q\log^3q)$ versus $\mathcal{O}(q^{2/3}\log q)$. Extra optimization/noise logarithms and larger leading constant. \\
\midrule
\textbf{(10)} Clipped SGD \citep{sgacGaash2025v2}, Thm.~1; derived bound on $\mathbb R^d$, sub-Gaussian noise: $\mathbb E\exp(\|g-\nabla f\|^2/\sigma^2)\le e$.
& $\mathcal{O}\!\left((1+T+q^2)H+v^2H^3\right)$ stochastic-gradient calls, including double sampling.
& \textbf{(6)}: $T\to\sqrt T$ and $q^2\to q$ up to logs on the sub-Gaussian overlap; extends to finite variance. Unchanged noise can mask a total gain.
& \textbf{(6)}: extra logarithmic factors on optimization terms. \\
\midrule
\textbf{(11)} \citep{fvRadiusRudomanov2026v1}, Thm.~3.7; special case $\sigma=0$, $L_1>0$, $\mathbb R^d$.
& $\mathcal{O}\!\left(\sqrt T+(1+q)\log(2+q)\right)$ exact-gradient calls, including line-search queries.
& \textbf{(6)}: extends the radius-only guarantee to finite-variance stochastic gradients; retains linear curvature dependence up to logarithms.
& \textbf{(6)}: at $\sigma=0$, pays $H^3$ on $\sqrt T$ and $(1+q)$, versus no logarithm and $\log(2+q)$, respectively. \\
\bottomrule
\end{tabularx}
\begin{tablenotes}[flushleft]\footnotesize
\item[] \textbf{Abbr:} $T=L_0R^2/\varepsilon$, $q=L_1R$, $v=\sigma R/\varepsilon$, $\ell=\log(e/\delta)$, $\lambda=1+\log(1+\Delta/\varepsilon)$, $W=1+\sqrt T+\lambda+(q\lambda)^{2/3}$, $M=1+T+q+\lambda+v^2\ell^2$; $J=1+\min\{\log(\Delta/\varepsilon),\log_+(L_1^2\Delta/L_0)\}$ in \textbf{(3)}. $\log_+u=\max\{0,\log u\}$, with $\log_+0=0$. The logarithmic fraction in \textbf{(1)} is zero at $v=0$. $c>0$ is fixed; constants in $\mathcal{O}_c$ may depend on $c$. $p$ in prior-work rows denotes the noise moment order.
\end{tablenotes}
\end{threeparttable}
\end{table}
\clearpage
